当然。FFT 和 NTT 都是用来**快速做多项式乘法/卷积**的算法家族。它们的核心思想相似：把多项式从“系数表示”转换到“点值表示”，逐点相乘，再转换回来。

**1. FFT 是什么**

FFT，全称 Fast Fourier Transform，快速傅里叶变换。

它用于快速计算离散傅里叶变换 DFT：

```text
A_k = sum_{j=0}^{n-1} a_j * omega_n^{jk}
```

其中 `omega_n` 是复数单位根，比如：

```text
omega_n = e^{2πi/n}
```

普通 DFT 需要 `O(n^2)`，FFT 利用单位根的对称性，把复杂度降到：

```text
O(n log n)
```

在算法竞赛和工程里，FFT 常用于：

- 大整数乘法
- 多项式乘法
- 信号处理
- 图像处理
- 卷积计算

例如两个多项式：

```text
A(x) = a0 + a1x + a2x^2
B(x) = b0 + b1x + b2x^2
```

朴素乘法是两层循环，复杂度 `O(n^2)`；用 FFT 可以做到 `O(n log n)`。

**2. NTT 是什么**

NTT，全称 Number Theoretic Transform，数论变换。

它可以理解为“在模意义下的 FFT”。FFT 使用复数单位根，而 NTT 使用模数域里的原根。

NTT 的形式类似：

```text
A_k = sum_{j=0}^{n-1} a_j * g^{jk} mod p
```

其中：

- `p` 是质数模数
- `g` 是模 `p` 意义下的合适单位根
- 所有运算都在整数取模下进行

NTT 的优势是：

- 没有浮点误差
- 结果精确
- 很适合算法竞赛和密码学相关计算

常见 NTT 模数：

```text
998244353 = 119 * 2^23 + 1
```

它很常用，因为它支持长度最高到 `2^23` 的 NTT，并且原根通常取 `3`。

**3. FFT 和 NTT 的区别**

| 对比项 | FFT | NTT |
|---|---|---|
| 运算对象 | 复数 | 模整数 |
| 是否有误差 | 有浮点误差 | 精确 |
| 常用场景 | 信号处理、大数乘法 | 竞赛、多项式、模运算 |
| 单位根 | 复数单位根 | 模意义下的原根 |
| 模数限制 | 无 | 需要特殊质数 |
| 实现难度 | 要处理精度 | 要处理模数和原根 |

**4. 它们为什么能加速乘法**

多项式乘法本质上是卷积：

```text
c_k = sum a_i * b_j    where i + j = k
```

直接算每个 `c_k` 是 `O(n^2)`。

FFT/NTT 的套路是：

```text
系数表示 -> 点值表示 -> 点值逐项相乘 -> 系数表示
```

也就是：

```text
A, B
=> FFT(A), FFT(B)
=> FFT(A) * FFT(B)
=> inverse FFT
=> C
```

NTT 同理，只是所有步骤都在模数下做。

**一句话总结**

FFT 是用复数做快速卷积，速度快但可能有精度误差；NTT 是用模整数做快速卷积，结果精确，但需要合适的质数模数。两者思想几乎一致，都是把多项式乘法从 `O(n^2)` 优化到 `O(n log n)`。